% Unofficial New York University Poster Template
% a fork of https://github.com/anishathalye/gemini
% also refer to https://github.com/k4rtik/uchicago-poster



\documentclass[final]{beamer}

% ====================
% Packages
% ====================

\usepackage[T1]{fontenc}
 \usepackage[utf8]{luainputenc}
\usepackage{lmodern}
\usepackage[size=custom, width=140,height=70, scale=1]{beamerposter}
\usetheme{gemini}
\usecolortheme{cam}
\usepackage{graphicx}
\usepackage{booktabs}
\usepackage{tikz}
\usepackage{pgfplots}
\pgfplotsset{compat=1.14}
\usepackage{anyfontsize}

% ====================
% Lengths
% ====================

% If you have N columns, choose \sepwidth and \colwidth such that
% (N+1)*\sepwidth + N*\colwidth = \paperwidth
\newlength{\sepwidth}
\newlength{\colwidth}
\setlength{\sepwidth}{0.025\paperwidth}
\setlength{\colwidth}{0.3\paperwidth}

\newcommand{\separatorcolumn}{\begin{column}{\sepwidth}\end{column}}

% ====================
% Title
% ====================

\title{Some fancy title: followed by some more text}

\author{M.Eng. Aoping Wu}

\institute[shortinst]{School of Transportation and Logistics, Southwest Jiaotong University}

% ====================
% Footer (optional)
% ====================

\footercontent{
  \href{https://github.com}{Github: https://github.com} \hfill
  TRB Conference 2023, Washington D.C. \hfill
  \href{mailto:wuaoping@my.swjtu.edu.cn}{wuaoping@my.swjtu.edu.cn}}
% (can be left out to remove footer)

% ====================
% Logo (optional)
% ====================

% use this to include logos on the left and/or right side of the header:
% Left: institution
 \logoleft{\includegraphics[height=5cm]{logos/swjtu_loge.jpg}}
% Right: funding agencies and other affilations 
%\logoright{\includegraphics[height=7cm]{logos/NSF.eps}}
% ====================
% Body
% ====================

\logoright{\includegraphics[height=5cm]{logos/school_logo.jpg}}

\begin{document}



\begin{frame}[t]
\begin{columns}[t]
\separatorcolumn

\begin{column}{\colwidth}

\begin{block}{ABOUT THIS TEMPLATE}

    This is a LaTeX follower from Southwest Jiaotong University.
    In the following content, I humbly introduce some examples usually used in the LaTeX, and the guidance for the TRB conference posters.
    I hope this template can help someone who is a new fish for the LaTeX and working for TRB posters, especially the fellow from SWJTU.

  \end{block}
  
  \begin{block}{The Poster sessions}

    Poster sessions are a valuable opportunity for authors to present papers and meet with interested attendees for in-depth technical discussions. TRB poster sessions are very well attended and a popular feature of the Annual Meeting. Good planning can make your presentation clear, effective, and rewarding.

  \end{block}

  \begin{block}{The goals in designing a poster}

    The items can be arranged as follows: 
    
    \begin{itemize}
      \item \textbf{Point 1} To attract attention.
      \item \textbf{Point 2} To provide a clear overview of your work.
      \item \textbf{Point 3} To provide enough material to explain the research without an oral explanation.
      \item \textbf{Point 4} To provide enough material to initiate discussion and questions without overwhelming the audience.
    \end{itemize}


  \end{block}

  \begin{alertblock}{What TRB Provides}
   

    Poster sessions are scheduled for 1 hour and 45 minutes in a large hall at the Convention Center. Several poster sessions will be taking place at once, with hundreds of papers being presented in each session time slot.

    TRB provides a display board that is four feet tall (4 ft. or 122 cm.) and eight feet wide (8 ft. or 244 cm.). A narrow table (6ft. long and 10 in. wide) is provided in front of the display board that may be used for material or a laptop computer.
    
    TRB provides thumbtacks for presenters to use to attach posters made of paper or lightweight poster board. Each poster display board is numbered for easy identification by both authors and attendees.


    \textbf{The dimensions of this poster have been revised according to the rules of The TRB conference.}

    \textit{No electrical outlets will be provided. Please come with your computer charged!}

  \end{alertblock}

\begin{block}{What Authors Should Bring?}

    \begin{enumerate}
      \item A single sheet poster made of heavyweight poster paper or similar material, approximately 4 ft. X 8 ft. in size.
      \item Heavy duty fasteners or tacks as back up to the thumbtacks. Bring tape or other materials as determined necessary.
      \item Authors may bring additional written material of a non-commercial nature that supplements the material presented on the poster. 
      \item Business cards. This is a GREAT networking opportunity. 
      \item \textbf{Comfortable shoes. Two hours on your feet can be tiring!}
    \end{enumerate}

  \end{block}

\end{column}

\separatorcolumn

\begin{column}{\colwidth}

  

  \begin{block}{How to insert picture?}

    Like these lines:

    \begin{figure}
        \centering
        \includegraphics[width=0.2\textwidth]{logos/1293741.jpg}
        \caption{The Wednesday Adans}
        \label{fig: The Wednesday Adans}
    \end{figure}

You can also cite this Figure \ref{fig: The Wednesday Adans}.
  \end{block}

  \begin{block}{Poster Session Rules}

\begin{itemize}

   \item Computers without audio may be used to run demonstrations and display additional information or illustrations. Telephone connections and other equipment are not allowed.
\item Commercial advertising of products or services is not permitted.
\item Internet will be available at the conference center, but presenters should be prepared in case the reception is poor due to overuse.
\item There are no provisions for making posters at the meeting, nor for receiving, storing, or returning posters to authors.
\item Poster display boards are rented and cannot be written on or defaced.
\end{itemize}

   

  \end{block}

\begin{block}{Poster Design Guidelines}
    \begin{enumerate}
        \item Prepare poster on a SINGLE sheet of heavy duty paper or similarly flexible material. The poster should be about 4 ft. X 8 ft. in size. Do not tack individual pages of a PowerPoint presentation or a text manuscript onto a poster
board.
\item Keep content SIMPLE. A poster is a visual communication tool, not a manuscript. The viewer should be able to easily identify the primary concepts of the project without wading through a lot of text or complex formulas. Identify 3 or 4 main points or concepts to communicate.
\item Present text in bullets or small chunks broken up by subheadings. Use at least 28–36 point bold sans serif font (e.g., Arial or Helvetica) for headers and 18–24 point font for text.
\item Present information in columns. Arrange material in a logical sequence, from left top to bottom right. Three columns is a good target to shoot for.
\item Offer a balanced mix of text and graphics. Too many words will result in people glossing over or simply bypassing your poster. A good rule of thumb is 50\% text, 50\% graphics and photos.
\item Avoid acronyms and jargon. Simple language is best.
\item Avoid dark-colored backgrounds. Use light colored backgrounds with black or very dark colored text. Graphics should similarly provide a stark contrast to be readable.
\item Use simple graphics. Charts, drawings and illustrations should be limited to a 2-3 color palette at a resolution of at least 300 dpi. Visuals should be large enough to be comfortably read from 3 feet away.
\item Provide author name(s), organization logos and/or other acknowledgements to give credit to those who have done the work.
\item Prepare a brief (up to 5 minutes) oral presentation for delivering to small audiences gathered around the poster.
    \end{enumerate}
\end{block}





\end{column}

\separatorcolumn

\begin{column}{\colwidth}

  \begin{exampleblock}{How to insert some math formulations?}

    A different kind of highlighted block.

    $$
    \int_{-\infty}^{\infty} e^{-x^2}\,dx = \sqrt{\pi}
    $$

    \heading{A heading inside a block}

    Praesent consectetur mi $x^2 + y^2$ metus, nec vestibulum justo viverra
    nec. Proin eget nulla pretium, egestas magna aliquam, mollis neque. Vivamus
    dictum $\mathbf{u}^\intercal\mathbf{v}$ sagittis odio, vel porta erat
    congue sed. Maecenas ut dolor quis arcu auctor porttitor.

    \heading{Another heading inside a block}

    Sed augue erat, scelerisque a purus ultricies, placerat porttitor neque.
    Donec $P(y \mid x)$ fermentum consectetur $\nabla_x P(y \mid x)$ sapien
    sagittis egestas. Duis eget leo euismod nunc viverra imperdiet nec id
    justo.

  \end{exampleblock}

  \begin{block}{How to insert tables?}

    Like this:

    \begin{table}
      \centering
      \begin{tabular}{l r r c}
        \toprule
        \textbf{First column} & \textbf{Second column} & \textbf{Third column} & \textbf{Fourth} \\
        \midrule
        Foo & 13.37 & 384,394 & $\alpha$ \\
        Bar & 2.17 & 1,392 & $\beta$ \\
        Baz & 3.14 & 83,742 & $\delta$ \\
        Qux & 7.59 & 974 & $\gamma$ \\
        \bottomrule
      \end{tabular}
      \caption{A table caption.}
    \end{table}



  \end{block}

  \begin{block}{References}
List the references with the \textbf{BibTeX} format in the file named \textbf{poster.bib}.
    \nocite{*}
    \footnotesize{\bibliographystyle{plain}\bibliography{poster}}

  \end{block}

\end{column}

\separatorcolumn
\end{columns}
\end{frame}

\end{document}
